%=====================================
%   20250528.tex
%   20250528 
%   toshi2019
%=====================================
% ------------------------
% documentclass
% ------------------------
\documentclass[leqno, 10pt, a4paper, dvipdfmx]{jsarticle}
\title{ゼミノート（05月28日分）}
\author{Toshi2019}
\date{\today}

% ------------------------
% usepackage
% ------------------------
\usepackage{algorithm}
\usepackage{algorithmic}
\usepackage{amscd}
\usepackage{amsfonts}
\usepackage{amsmath}
\usepackage[psamsfonts]{amssymb}
\usepackage{amsthm}
\usepackage{ascmac}
\usepackage{bm}
\usepackage{color}
\usepackage{enumerate}
\usepackage{fancybox}
\usepackage[stable]{footmisc}
\usepackage{graphicx}
\usepackage{listings}
\usepackage{mathrsfs}
\usepackage{mathtools}
\usepackage{otf}
\usepackage{pifont}
\usepackage{proof}
\usepackage{subfigure}
\usepackage{tikz}
\usepackage{verbatim}
\usepackage[all]{xy}

\usetikzlibrary{cd}
\usetikzlibrary{arrows.meta}


% ================================
% パッケージを追加する場合のスペース 
\usepackage[dvipdfmx]{hyperref}
\usepackage{xcolor}
\definecolor{darkgreen}{rgb}{0,0.45,0} 
\definecolor{darkred}{rgb}{0.75,0,0}
\definecolor{darkblue}{rgb}{0,0,0.6} 
\hypersetup{
    colorlinks=true,
    citecolor=darkgreen,
    linkcolor=darkred,
    urlcolor=darkblue,
}
\usepackage{pxjahyper}

%\usepackage[mathscr]{euscript} % mathscr のフォントを変更

%\usepackage{mmasym}

%=================================


% --------------------------
% theoremstyle
% --------------------------
\theoremstyle{definition}

% --------------------------
% newtheoem
% --------------------------

% 日本語で定理, 命題, 証明などを番号付きで用いるためのコマンドです. 
% If you want to use theorem environment in Japanece, 
% you can use these code. 
% Attention!
% All theorem enivironment numbers depend on 
% only section numbers.
\newtheorem{Axiom}{公理}[section]
\newtheorem{Definition}[Axiom]{定義}
\newtheorem{Theorem}[Axiom]{定理}
\newtheorem{Proposition}[Axiom]{命題}
\newtheorem{Lemma}[Axiom]{補題}
\newtheorem{Corollary}[Axiom]{系}
\newtheorem{Example}[Axiom]{例}
\newtheorem{Claim}[Axiom]{主張}
\newtheorem{Property}[Axiom]{性質}
\newtheorem{Attention}[Axiom]{注意}
\newtheorem{Question}[Axiom]{問}
\newtheorem{Problem}[Axiom]{問題}
\newtheorem{Consideration}[Axiom]{考察}
\newtheorem{Alert}[Axiom]{警告}
\newtheorem{Fact}[Axiom]{事実}


% 日本語で定理, 命題, 証明などを番号なしで用いるためのコマンドです. 
% If you want to use theorem environment with no number in Japanese, You can use these code.
\newtheorem*{Axiom*}{公理}
\newtheorem*{Definition*}{定義}
\newtheorem*{Theorem*}{定理}
\newtheorem*{Proposition*}{命題}
\newtheorem*{Lemma*}{補題}
\newtheorem*{Example*}{例}
\newtheorem*{Corollary*}{系}
\newtheorem*{Claim*}{主張}
\newtheorem*{Property*}{性質}
\newtheorem*{Attention*}{注意}
\newtheorem*{Question*}{問}
\newtheorem*{Problem*}{問題}
\newtheorem*{Consideration*}{考察}
\newtheorem*{Alert*}{警告}
\newtheorem{Fact*}{事実}


% 英語で定理, 命題, 証明などを番号付きで用いるためのコマンドです. 
%　If you want to use theorem environment in English, You can use these code.
%all theorem enivironment number depend on only section number.
\newtheorem{Axiom+}{Axiom}[section]
\newtheorem{Definition+}[Axiom+]{Definition}
\newtheorem{Theorem+}[Axiom+]{Theorem}
\newtheorem{Proposition+}[Axiom+]{Proposition}
\newtheorem{Lemma+}[Axiom+]{Lemma}
\newtheorem{Example+}[Axiom+]{Example}
\newtheorem{Corollary+}[Axiom+]{Corollary}
\newtheorem{Claim+}[Axiom+]{Claim}
\newtheorem{Property+}[Axiom+]{Property}
\newtheorem{Attention+}[Axiom+]{Attention}
\newtheorem{Question+}[Axiom+]{Question}
\newtheorem{Problem+}[Axiom+]{Problem}
\newtheorem{Consideration+}[Axiom+]{Consideration}
\newtheorem{Alert+}{Alert}
\newtheorem{Fact+}[Axiom+]{Fact}
\newtheorem{Remark+}[Axiom+]{Remark}

% ----------------------------
% commmand
% ----------------------------
% 執筆に便利なコマンド集です. 
% コマンドを追加する場合は下のスペースへ. 

% 集合の記号 (黒板文字)
\newcommand{\NN}{\mathbb{N}}
\newcommand{\ZZ}{\mathbb{Z}}
\newcommand{\QQ}{\mathbb{Q}}
\newcommand{\RR}{\mathbb{R}}
\newcommand{\CC}{\mathbb{C}}
\newcommand{\PP}{\mathbb{P}}
\newcommand{\KK}{\mathbb{K}}


% 集合の記号 (太文字)
\newcommand{\nn}{\mathbf{N}}
\newcommand{\zz}{\mathbf{Z}}
\newcommand{\qq}{\mathbf{Q}}
\newcommand{\rr}{\mathbf{R}}
\newcommand{\cc}{\mathbf{C}}
\newcommand{\pp}{\mathbf{P}}
\newcommand{\kk}{\mathbf{k}}

% 特殊な写像の記号
\newcommand{\ev}{\mathop{\mathrm{ev}}\nolimits} % 値写像
\newcommand{\pr}{\mathop{\mathrm{pr}}\nolimits} % 射影

% スクリプト体にするコマンド
%   例えば　{\mcal C} のように用いる
\newcommand{\mcal}{\mathcal}

% 花文字にするコマンド 
%   例えば　{\h C} のように用いる
\newcommand{\h}{\mathscr}

% ヒルベルト空間などの記号
\newcommand{\F}{\mcal{F}}
\newcommand{\X}{\mcal{X}}
\newcommand{\Y}{\mcal{Y}}
\newcommand{\Hil}{\mcal{H}}
\newcommand{\RKHS}{\Hil_{k}}
\newcommand{\Loss}{\mcal{L}_{D}}
\newcommand{\MLsp}{(\X, \Y, D, \Hil, \Loss)}

% 偏微分作用素の記号
\newcommand{\p}{\partial}

% 角カッコの記号 (内積は下にマクロがあります)
\newcommand{\lan}{\langle}
\newcommand{\ran}{\rangle}



% 圏の記号など
\newcommand{\Set}{{\bf Set}}
\newcommand{\Vect}{{\bf Vect}}
\newcommand{\FDVect}{{\bf FDVect}}
%\newcommand{\Ring}{{\bf Ring}}
\newcommand{\Ab}{{\bf Ab}}
\newcommand{\Mod}{\mathop{\mathrm{Mod}}\nolimits}
\newcommand{\Modf}{\mathop{\mathrm{Mod}^\mathrm{f}}\nolimits}
\newcommand{\CGA}{{\bf CGA}}
\newcommand{\GVect}{{\bf GVect}}
\newcommand{\Lie}{{\bf Lie}}
\newcommand{\dLie}{{\bf Liec}}



% 射の集合など
\newcommand{\Map}{\mathop{\mathrm{Map}}\nolimits} % 写像の集合
\newcommand{\Hom}{\mathop{\mathrm{Hom}}\nolimits} % 射集合
\newcommand{\End}{\mathop{\mathrm{End}}\nolimits} % 自己準同型の集合
\newcommand{\Aut}{\mathop{\mathrm{Aut}}\nolimits} % 自己同型の集合
\newcommand{\Mor}{\mathop{\mathrm{Mor}}\nolimits} % 射集合
\newcommand{\Ker}{\mathop{\mathrm{Ker}}\nolimits} % 核
\newcommand{\Img}{\mathop{\mathrm{Im}}\nolimits} % 像
\newcommand{\Cok}{\mathop{\mathrm{Coker}}\nolimits} % 余核
\newcommand{\Cim}{\mathop{\mathrm{Coim}}\nolimits} % 余像

% その他便利なコマンド
\newcommand{\dip}{\displaystyle} % 本文中で数式モード
\newcommand{\e}{\varepsilon} % イプシロン
\newcommand{\dl}{\delta} % デルタ
\newcommand{\pphi}{\varphi} % ファイ
\newcommand{\ti}{\tilde} % チルダ
\newcommand{\pal}{\parallel} % 平行
\newcommand{\op}{{\rm op}} % 双対を取る記号
\newcommand{\lcm}{\mathop{\mathrm{lcm}}\nolimits} % 最小公倍数の記号
\newcommand{\Probsp}{(\Omega, \F, \P)} 
\newcommand{\argmax}{\mathop{\rm arg~max}\limits}
\newcommand{\argmin}{\mathop{\rm arg~min}\limits}





% ================================
% コマンドを追加する場合のスペース 
%\newcommand{\OO}{\mcal{O}}



\makeatletter
\renewenvironment{proof}[1][\proofname]{\par
  \pushQED{\qed}%
  \normalfont \topsep6\p@\@plus6\p@\relax
  \trivlist
  \item[\hskip\labelsep
%        \itshape
         \bfseries
%    #1\@addpunct{.}]\ignorespaces
    {#1}]\ignorespaces
}{%
  \popQED\endtrivlist\@endpefalse
}
\makeatother

\renewcommand{\proofname}{証明.}



%\renewcommand\proofname{\bf 証明} % 証明
\numberwithin{equation}{section}
\newcommand{\cTop}{\textsf{Top}}
%\newcommand{\cOpen}{\textsf{Open}}
\newcommand{\Op}{\mathop{\textsf{Op}}\nolimits}
\newcommand{\Ob}{\mathop{\textrm{Ob}}\nolimits}
\newcommand{\id}{\mathop{\mathrm{id}}\nolimits}
\newcommand{\pt}{\mathop{\mathrm{pt}}\nolimits}
\newcommand{\res}{\mathop{\rho}\nolimits}
\newcommand{\A}{\mcal{A}}
\newcommand{\B}{\mcal{B}}
\newcommand{\C}{\mcal{C}}
\newcommand{\D}{\mcal{D}}
\newcommand{\E}{\mcal{E}}
\newcommand{\G}{\mcal{G}}
%\newcommand{\H}{\mcal{H}}
\newcommand{\I}{\mcal{I}}
\newcommand{\J}{\mcal{J}}
\newcommand{\OO}{\mcal{O}}
\newcommand{\Ring}{\mathop{\textsf{Ring}}\nolimits}
\newcommand{\cAb}{\mathop{\textsf{Ab}}\nolimits}
%\newcommand{\Ker}{\mathop{\mathrm{Ker}}\nolimits}
\newcommand{\im}{\mathop{\mathrm{Im}}\nolimits}
\newcommand{\Coker}{\mathop{\mathrm{Coker}}\nolimits}
\newcommand{\Coim}{\mathop{\mathrm{Coim}}\nolimits}
\newcommand{\rank}{\mathop{\mathrm{rank}}\nolimits}
\newcommand{\Ht}{\mathop{\mathrm{Ht}}\nolimits}
\newcommand{\supp}{\mathop{\mathrm{supp}}\nolimits}
\newcommand{\colim}{\mathop{\mathrm{colim}}}
\newcommand{\Tor}{\mathop{\mathrm{Tor}}\nolimits}

\newcommand{\cat}{\mathscr{C}}

%筆記体
\newcommand{\cA}{\mcal{A}}
\newcommand{\cB}{\mcal{B}}
\newcommand{\cC}{\mcal{C}}
\newcommand{\cD}{\mcal{D}}
\newcommand{\cE}{\mcal{E}}
\newcommand{\cF}{\mcal{F}}
\newcommand{\cG}{\mcal{G}}
\newcommand{\cH}{\mcal{H}}
\newcommand{\cI}{\mcal{I}}
\newcommand{\cJ}{\mcal{J}}
\newcommand{\cK}{\mcal{K}}
\newcommand{\cL}{\mcal{L}}
\newcommand{\cM}{\mcal{M}}
\newcommand{\cN}{\mcal{N}}
\newcommand{\cO}{\mcal{O}}
\newcommand{\cP}{\mcal{P}}
\newcommand{\cQ}{\mcal{Q}}
\newcommand{\cR}{\mcal{R}}
\newcommand{\cS}{\mcal{S}}
\newcommand{\cT}{\mcal{T}}
\newcommand{\cU}{\mcal{U}}
\newcommand{\cV}{\mcal{V}}
\newcommand{\cW}{\mcal{W}}
\newcommand{\cX}{\mcal{X}}
\newcommand{\cY}{\mcal{Y}}
\newcommand{\cZ}{\mcal{Z}}


\newcommand{\scA}{\mathscr{A}}
\newcommand{\scB}{\mathscr{B}}
\newcommand{\scC}{\mathscr{C}}
\newcommand{\scD}{\mathscr{D}}
\newcommand{\scE}{\mathscr{E}}
\newcommand{\scF}{\mathscr{F}}
\newcommand{\scN}{\mathscr{N}}
\newcommand{\scO}{\mathscr{O}}
\newcommand{\scV}{\mathscr{V}}
\newcommand{\scU}{\mathscr{U}}


\newcommand{\ibA}{\mathop{\text{\textit{\textbf{A}}}}}
\newcommand{\ibB}{\mathop{\text{\textit{\textbf{B}}}}}
\newcommand{\ibC}{\mathop{\text{\textit{\textbf{C}}}}}
\newcommand{\ibD}{\mathop{\text{\textit{\textbf{D}}}}}
\newcommand{\ibE}{\mathop{\text{\textit{\textbf{E}}}}}
\newcommand{\ibF}{\mathop{\text{\textit{\textbf{F}}}}}
\newcommand{\ibG}{\mathop{\text{\textit{\textbf{G}}}}}
\newcommand{\ibH}{\mathop{\text{\textit{\textbf{H}}}}}
\newcommand{\ibI}{\mathop{\text{\textit{\textbf{I}}}}}
\newcommand{\ibJ}{\mathop{\text{\textit{\textbf{J}}}}}
\newcommand{\ibK}{\mathop{\text{\textit{\textbf{K}}}}}
\newcommand{\ibL}{\mathop{\text{\textit{\textbf{L}}}}}
\newcommand{\ibM}{\mathop{\text{\textit{\textbf{M}}}}}
\newcommand{\ibN}{\mathop{\text{\textit{\textbf{N}}}}}
\newcommand{\ibO}{\mathop{\text{\textit{\textbf{O}}}}}
\newcommand{\ibP}{\mathop{\text{\textit{\textbf{P}}}}}
\newcommand{\ibQ}{\mathop{\text{\textit{\textbf{Q}}}}}
\newcommand{\ibR}{\mathop{\text{\textit{\textbf{R}}}}}
\newcommand{\ibS}{\mathop{\text{\textit{\textbf{S}}}}}
\newcommand{\ibT}{\mathop{\text{\textit{\textbf{T}}}}}
\newcommand{\ibU}{\mathop{\text{\textit{\textbf{U}}}}}
\newcommand{\ibV}{\mathop{\text{\textit{\textbf{V}}}}}
\newcommand{\ibW}{\mathop{\text{\textit{\textbf{W}}}}}
\newcommand{\ibX}{\mathop{\text{\textit{\textbf{X}}}}}
\newcommand{\ibY}{\mathop{\text{\textit{\textbf{Y}}}}}
\newcommand{\ibZ}{\mathop{\text{\textit{\textbf{Z}}}}}

\newcommand{\ibx}{\mathop{\text{\textit{\textbf{x}}}}}

%\newcommand{\Comp}{\mathop{\mathrm{C}}\nolimits}
%\newcommand{\Komp}{\mathop{\mathrm{K}}\nolimits}
%\newcommand{\Domp}{\mathop{\mathsf{D}}\nolimits}%複体のホモトピー圏
%\newcommand{\Comp}{\mathrm{C}}
%\newcommand{\Komp}{\mathrm{K}}
%\newcommand{\Domp}{\mathsf{D}}%複体のホモトピー圏

\newcommand{\Comp}{\mathop{\mathrm{C}}\nolimits}
\newcommand{\Komp}{\mathop{\mathsf{K}}\nolimits}
\newcommand{\Domp}{\mathord{\mathsf{D}}}%\nolimits}
\newcommand{\Kompl}{\mathop{\mathsf{K}^\mathrm{+}}\nolimits}
\newcommand{\Kompu}{\mathop{\mathsf{K}^\mathrm{-}}\nolimits}
\newcommand{\Kompb}{\mathop{\mathsf{K}^\mathrm{b}}\nolimits}
\newcommand{\Dompl}{\mathop{\mathsf{D}^\mathrm{+}}\nolimits}
\newcommand{\Dompu}{\mathop{\mathsf{D}^\mathrm{-}}\nolimits}
\newcommand{\Dompb}{\mathop{\mathsf{D}^\mathrm{b}}\nolimits}
\newcommand{\Dbconic}{\mathop{\mathsf{D}^\mathrm{b}_{\rrp}}\nolimits}
\newcommand{\Dompbf}{\mathop{\mathsf{D}_\mathrm{f}^\mathrm{b}}\nolimits}
\newcommand{\Domplb}{\mathop{\mathsf{D}^\mathrm{lb}}\nolimits}




\newcommand{\CCat}{\Comp(\cat)}
\newcommand{\KCat}{\Komp(\cat)}
\newcommand{\DCat}{\Domp(\cat)}%圏Cの複体のホモトピー圏
\newcommand{\HOM}{\mathop{\mathscr{H}\hspace{-2pt}om}\nolimits}%内部Hom
\newcommand{\RHOM}{\mathop{\mathrm{R}\hspace{-1.5pt}\HOM}\nolimits}

\newcommand{\muS}{\mathop{\mathrm{SS}}\nolimits}
\newcommand{\RG}{\mathop{\mathrm{R}\hspace{-0pt}\Gamma}\nolimits}
\newcommand{\RHom}{\mathop{\mathrm{R}\hspace{-1.5pt}\Hom}\nolimits}
\newcommand{\Rder}{\mathrm{R}}

\newcommand{\simar}{\mathrel{\overset{\sim}{\rightarrow}}}%同型右矢印
\newcommand{\simarr}{\mathrel{\overset{\sim}{\longrightarrow}}}%同型右矢印
\newcommand{\simra}{\mathrel{\overset{\sim}{\leftarrow}}}%同型左矢印
\newcommand{\simrra}{\mathrel{\overset{\sim}{\longleftarrow}}}%同型左矢印

\newcommand{\hocolim}{{\mathrm{hocolim}}}
\newcommand{\indlim}[1][]{\mathop{\varinjlim}\limits_{#1}}
\newcommand{\sindlim}[1][]{\smash{\mathop{\varinjlim}\limits_{#1}}\,}
\newcommand{\Pro}{\mathrm{Pro}}
\newcommand{\Ind}{\mathrm{Ind}}
\newcommand{\prolim}[1][]{\mathop{\varprojlim}\limits_{#1}}
\newcommand{\sprolim}[1][]{\smash{\mathop{\varprojlim}\limits_{#1}}\,}

\newcommand{\Sh}{\mathrm{Sh}}
\newcommand{\PSh}{\mathrm{PSh}}

\newcommand{\rmD}{\mathsf{D}}
\newcommand{\vD}{\mathbf{D}}

\newcommand{\Lloc}[1][]{\mathord{\mathcal{L}^1_{\mathrm{loc},{#1}}}}
\newcommand{\ori}{\mathord{\mathrm{or}}}
\newcommand{\Db}{\mathord{\mathscr{D}b}}

\newcommand{\codim}{\mathop{\mathrm{codim}}\nolimits}



\newcommand{\gld}{\mathop{\mathrm{gld}}\nolimits}
\newcommand{\wgld}{\mathop{\mathrm{wgld}}\nolimits}


\newcommand{\tens}[1][]{\mathbin{\otimes_{\raise1.5ex\hbox to-.1em{}{#1}}}}
\newcommand{\etens}{\mathbin{\boxtimes}}
\newcommand{\ltens}[1][]{\mathbin{\overset{\mathrm{L}}\tens}_{#1}}
\newcommand{\mtens}[1][]{\mathbin{\overset{\mathrm{\mu}}\tens}_{#1}}
\newcommand{\lltens}[1][]{\mathbin{{\mathop{\tens}\limits^{\mathrm{L}}_{#1}}}}
\newcommand{\letens}{\overset{\mathrm{L}}{\etens}}
\newcommand{\detens}{\underline{\etens}}
\newcommand{\ldetens}{\overset{\mathrm{L}}{\underline{\etens}}}
\newcommand{\dtens}[1][]{{\overset{\mathrm{L}}{\underline{\otimes}}}_{#1}}

\newcommand{\blk}{\mathord{\ \cdot\ }}
\newcommand{\mres}[2][]{{\left.{#1}\right\rvert}_{#2}}


%\newcommand{\hocolim}{{\mathrm{hocolim}}}
%\newcommand{\indlim}[1][]{\mathop{\varinjlim}\limits_{#1}}
%\newcommand{\sindlim}[1][]{\smash{\mathop{\varinjlim}\limits_{#1}}\,}
%\newcommand{\Pro}{\mathrm{Pro}}
%\newcommand{\Ind}{\mathrm{Ind}}
%\newcommand{\prolim}[1][]{\mathop{\varprojlim}\limits_{#1}}
%\newcommand{\sprolim}[1][]{\smash{\mathop{\varprojlim}\limits_{#1}}\,}
\newcommand{\proolim}[1][]{\mathop{\text{\rm``{$\varprojlim$}''}}\limits_{#1}}
\newcommand{\sproolim}[1][]{\smash{\mathop{\rm``{\varprojlim}''}\limits_{#1}}}
\newcommand{\inddlim}[1][]{\mathop{\text{\rm``{$\varinjlim$}''}}\limits_{#1}}
\newcommand{\sinddlim}[1][]{\smash{\mathop{\text{\rm``{$\varinjlim$}''}}\limits_{#1}}\,}
\newcommand{\ooplus}{\mathop{\text{\rm``{$\oplus$}''}}\limits}
\newcommand{\bbigsqcup}{\mathop{``\bigsqcup''}\limits}
\newcommand{\bsqcup}{\mathop{``\sqcup''}\limits}
\newcommand{\dsum}[1][]{\mathbin{\oplus_{#1}}}

\newcommand{\Fct}{\mathop{\mathsf{Fct}}\nolimits}

\newcommand{\pb}[1][]{\mathbin{\mathop{\times}\limits_{#1}}}
\newcommand{\dpb}[1][]{\mathbin{\mathop{\times}\limits_{#1}}}
\newcommand{\tpb}[1][]{\mathbin{\mathop{\times}\nolimits_{#1}}}





%================================================
% 自前の定理環境
%   https://mathlandscape.com/latex-amsthm/
% を参考にした
\newtheoremstyle{mystyle}%   % スタイル名
    {5pt}%                   % 上部スペース
    {5pt}%                   % 下部スペース
    {}%              % 本文フォント
    {}%                  % 1行目のインデント量
    {\bfseries}%                      % 見出しフォント
    {.}%                     % 見出し後の句読点
    {12pt}%                     % 見出し後のスペース
    {\thmname{#1}\thmnumber{ #2}\thmnote{{\hspace{2pt}\normalfont (#3)}}}% % 見出しの書式

\theoremstyle{mystyle}
\newtheorem{AXM}{公理}[section]
\newtheorem{DFN}[AXM]{定義}
\newtheorem{THM}[AXM]{定理}
\newtheorem*{THM*}{定理}
\newtheorem{PRP}[AXM]{命題}
\newtheorem{LMM}[AXM]{補題}
\newtheorem{CRL}[AXM]{系}
\newtheorem{EG}[AXM]{例}
\newtheorem*{EG*}{例}
\newtheorem{RMK}[AXM]{注意}
\newtheorem{CNV}[AXM]{約束}
\newtheorem{CMT}[AXM]{コメント}
\newtheorem*{CMT*}{コメント}
\newtheorem{NTN}[AXM]{記号}

% 定理環境ここまで
%====================================================

\usepackage{framed}
\definecolor{lightgray}{rgb}{0.75,0.75,0.75}
\renewenvironment{leftbar}{%
  \def\FrameCommand{\textcolor{lightgray}{\vrule width 4pt} \hspace{10pt}}% 
  \MakeFramed {\advance\hsize-\width \FrameRestore}}%
{\endMakeFramed}
\newenvironment{redleftbar}{%
  \def\FrameCommand{\textcolor{lightgray}{\vrule width 1pt} \hspace{10pt}}% 
  \MakeFramed {\advance\hsize-\width \FrameRestore}}%
 {\endMakeFramed}



\renewcommand{\Re}{\mathop{\mathrm{Re}}\nolimits}

% =================================





% ---------------------------
% new definition macro
% ---------------------------
% 便利なマクロ集です

% 内積のマクロ
%   例えば \inner<\pphi | \psi> のように用いる
\def\inner<#1>{\langle #1 \rangle}

% ================================
% マクロを追加する場合のスペース 

%=================================



% 位相空間
\newcommand{\Int}[1][]{\mathop{\mathrm{Int}}\nolimits_{#1}}
\newcommand{\Cl}[1][]{\mathop{\mathrm{Cl}}\nolimits_{#1}}
\newcommand{\Bd}[1][]{\mathop{\mathrm{Bd}}\nolimits_{#1}}
\newcommand{\bd}{\mathord{\partial}}
\newcommand{\Fr}[1][]{\mathop{\mathrm{Fr}}\nolimits_{#1}}


\newcommand{\mtan}[1][]{T\hspace{-1.75pt}{#1}}
\newcommand{\mtp}[1][]{{}^{t\!}{#1}} % transpose of the morphism
%\newcommand{\mpb}[1][]{{}^{t\!}{#1}'} % pull-back of the morphism
\newcommand{\mpb}[1][]{{#1}_{\pi}} % pull-back of the morphism
\newcommand{\mdiff}[1][]{{#1}_{d}} % pull-back of the morphism
\newcommand{\mptan}[1][]{{#1}_{\tau}} % pull-back of the morphism


%\newcommand{\pb}[1][]{\mathbin{\mathop{\times}\limits_{#1}}}

\newcommand{\cotb}[1][]{T^{\ast}{#1}}
\newcommand{\conm}[2][]{T_{#1}^{\hspace{0.7pt}\ast}{#2}}
\newcommand{\norb}[2][]{T_{#1}{#2}}

\newcommand{\cotz}[1][]{T^{\ast}_{#1}{#1}}

\newcommand{\cotn}[1][]{\dot{T}^{\ast}{#1}}


\newcommand{\muhom}{\mu hom}
\newcommand{\muHom}[1][]{\mathrm{Hom}^\mu_{\raise1.5ex\hbox to.1em{}#1}}

\newcommand{\mucat}[2][]{\mathsf{D}^{\mathrm{b}}_{#1}(#2)}
%\newcommand{\mucat}[2][]{\mathsf{D}^{\mathrm{b}}_{#1}(#2)}

\newcommand{\conv}[1][]{\mathbin{\mathop{\circ}\limits_{#1}}}

\newcommand{\torus}{\mathbf{T}}

\newcommand{\diag}[1][]{\Delta_{#1}}


% ----------------------------
% documenet 
% ----------------------------
% 以下, 本文の執筆スペースです. 
% Your main code must be written between 
% begin document and end document.
% ---------------------------

\begin{document}
\maketitle
\tableofcontents

\section{前回の同一視について}
前回，次の図式の射の構成について議論した．
\begin{equation}%\label{eq:4.4.1}
  \vcenter{\xymatrix@C=36pt@R=36pt{
    T^{\ast}_{\diag[Y]}(Y\times Y)
    \ar[d]^-{\rotatebox{90}{\(\sim\)}}
    &
    T^{\ast}_{\diag[f]}(X\times Y)
    \ar[d]^-{\rotatebox{90}{\(\sim\)}}
    \ar[l]
    \ar[r]
    &
    T^{\ast}_{\diag[X]}(X\times X)
    \ar[d]^-{\rotatebox{90}{\(\sim\)}}
    \\
    \cotb[Y]
    &
    Y\pb[X]\cotb[X]
    \ar[l]^-{\mdiff[f]}
    \ar[r]_-{\mpb[f]}
    &
    \cotb[X].
  }}\tag{4.4.1}
\end{equation}
これは，図式
\begin{equation}%\label{eq:4.4.3}
  \vcenter{\xymatrix@C=36pt@R=36pt{
    Y\times Y
    \ar[r]_-{f_1}
    &
    X\times Y
    \ar@{}[rd]|\square
    %\ar[d]^-{\rotatebox{90}{\(\sim\)}}
    \ar[r]_-{f_2}
    &
    X\times X
    %\ar[d]^-{\rotatebox{90}{\(\sim\)}}
    \\
    \diag[Y]  
    \ar@{^{(}->}[u]^-{}
    \ar[r]^-{{\sim}}
    &
    \diag[]
    \ar@{^{(}->}[u]^-{}
    %\ar[l]^-{\mdiff[f]}
    \ar[r]_-{f}
    &
    \diag[X].
    \ar@{^{(}->}[u]^-{}
  }}\tag{4.4.3}
\end{equation}
の\(f_2\)が\(\diag[X]\)に対して横断的であることを示す際に，
\(\mres[f_2]{\diag}\colon\diag[]\to\diag[X]\)の基本図式
\[
  \vcenter{\xymatrix@C=36pt@R=36pt{
  T^{\ast}_{\diag[]}(X\times Y)
  \ar[dr]^-{}
  &
  \diag[]\dpb[{\diag[X]}]T^{\ast}_{\diag[X]}(X\times X)
  \ar[d]^-{}
  \ar[l]_-{\mdiff[{f_2}_{\diag[]}]}
  \ar[r]^-{\mpb[{f_2}_{\diag[]}]}
  &
  T^{\ast}_{\diag[X]}(X\times X)
  \ar[d]^-{}
  \\
  %\cotb[Y]
  &
  \diag[]
  %\ar[l]^-{}
  \ar[r]_-{\mres[{f_2}]{\diag[]}}
  &
  \diag[X]
}}\]
の\(\mdiff[{f_2}_{\diag[]}]\)が同型であることから
従うことをコメントされた際の議論である．

\begin{proof}[{\(\mdiff[{f_2}_{\diag[]}]\)が同型であることの証明．}]
  単射であることは前回やった．
  したがって，\(\diag[]\)上のベクトル束の階数が
  一致することを言えばよい\footnote{
    これは，ベクトル束の射が底空間の点を保つことから，
    ファイバーごとに全射・単射を議論すれば良いことによる．
  }．
  \begin{align*}
    \rank{\conm[{\diag[]}]{(X\times Y)}}
    &=\codim{\diag[]}\\
    &=(\dim{X}+\dim{Y})-\dim{Y}\\
    &=\dim{X}
  \end{align*}であり，\begin{align*}
    \rank{\diag[]\dpb[{\diag[X]}]\conm[{\diag[X]}]{(X\times X)}}
    &=\rank{\conm[{\diag[X]}]{(X\times X)}}\quad \text{（引き戻しはファイバーを保つ）}\\
    &=\codim{\diag[]}\\
    &=(\dim{X}+\dim{X})-\dim{X}\\
    &=\dim{X}
  \end{align*}
  である．
\end{proof}
さらに，\(f_1\)の\(\diag[Y]\)への制限
\(\mres[f_1]{\diag[Y]}\colon\diag[Y]\to\diag[]\)の基本図式
を考える．
\[
  \vcenter{\xymatrix@C=36pt@R=36pt{
  T^{\ast}_{\diag[Y]}(Y\times Y)
  \ar[dr]^-{}
  &
  \diag[Y]\dpb[{\diag[]}]T^{\ast}_{\diag[]}(X\times Y)
  \ar[d]^-{}
  \ar[l]_-{\mdiff[{f_1}_{\diag[Y]}]}
  \ar[r]^-{\mpb[{f_1}_{\diag[Y]}]}_{\sim}
  &
  T^{\ast}_{\diag[]}(X\times Y)
  \ar[d]^-{}
  \\
  %\cotb[Y]
  &
  \diag[Y]
  %\ar[l]^-{}
  \ar[r]_-{\mres[{f_1}]{\diag[Y]}}^-{\sim}
  &
  \diag[]
}}\]
ここで，\(\mres[{f_1}]{{\diag[Y]}}\)は同型なので，
引き戻しの関手性から，\(\mpb[{f_1}_{\diag[Y]}]\)も同型である．


\section{関手\(\muhom\)}
%\newcommand{\mucat}[2][]{\mathsf{D}^{\mathrm{b}}_{#1}(#2)}
\(f\)を\(Y\)から\(X\)への射とする．
\(\diag[f]\subset X\times Y\)を\(f\)のグラフとする．
\(\diag[X]\)と\(\diag[Y]\)で
それぞれ\(X\times X\)と
\(Y\times Y\)の対角集合を表す．
\(Y\mathbin{\times_{X}}\cotb[X]\)と\(T^{\ast}_{\diag[f]}(X\times Y)\)を
射影\(\cotb[(X\times Y)]\to(\cotb[X])\times Y\)で同一視する．
同様に\(\cotb[X]\)と\(T^{\ast}_{\diag[X]}(X\times X)\)，
\(\cotb[Y]\)と\(T^{\ast}_{\diag[Y]}(Y\times Y)\)を
それぞれ第1射影で同一視する．
混同の恐れがないときは\(\diag[f]\)を\(\diag\)とかく．
すると，次の写像を得る（式 (4.3.3) 参照）．
\begin{equation}\label{eq:4.4.1}
  \vcenter{\xymatrix@C=36pt@R=36pt{
    T^{\ast}_{\diag[Y]}(Y\times Y)
    \ar[d]^-{\rotatebox{90}{\(\sim\)}}
    &
    T^{\ast}_{\diag[f]}(X\times Y)
    \ar[d]^-{\rotatebox{90}{\(\sim\)}}
    \ar[l]
    \ar[r]
    &
    T^{\ast}_{\diag[X]}(X\times X)
    \ar[d]^-{\rotatebox{90}{\(\sim\)}}
    \\
    \cotb[Y]
    &
    Y\pb[X]\cotb[X]
    \ar[l]^-{\mdiff[f]}
    \ar[r]_-{\mpb[f]}
    &
    \cotb[X].
  }}
\end{equation}
次の公式をよく用いる．
\begin{equation}\label{eq:4.4.2}
  \omega_{{Y\mathbin{\times_{X}}\cotb[X]}\mathbin{/}\cotb[Y]}
  \tens[]
  \omega_{{Y\mathbin{\times_{X}}\cotb[X]}\mathbin{/}\cotb[X]}
  \cong
  A_{Y\mathbin{\times_{X}}\cotb[X]}
\end{equation}
\begin{proof}
  \(
    \omega_{{Y\mathbin{\times_{X}}\cotb[X]}\mathbin{/}\cotb[Y]}
    \tens[]
    \omega_{{Y\mathbin{\times_{X}}\cotb[X]}\mathbin{/}\cotb[X]}
  \)は\(A_{Y\mathbin{\times_{X}}\cotb[X]}\)加群の複体なので射
  \[
    A_{Y\mathbin{\times_{X}}\cotb[X]}
    \to
    \omega_{{Y\mathbin{\times_{X}}\cotb[X]}\mathbin{/}\cotb[Y]}
    \tens[]
    \omega_{{Y\mathbin{\times_{X}}\cotb[X]}\mathbin{/}\cotb[X]}
  \]
  が存在する．したがって茎ごとに同型であることを示せばよい．
  \((y,f(x);\xi)\)を\(Y\mathbin{\times_{X}}\cotb[X]\)の点とする．
  \begin{align*}
    \left(
      \mdiff[f]^{!}A_{\cotb[Y]}\tens[]\mpb[f]^{!}A_{\cotb[X]}
    \right)_{(y,f(x);\xi)}
    &=\left(
      \mdiff[f]^{-1}A_{\cotb[Y]}\tens[]\mpb[f]^{-1}A_{\cotb[X]}
    \right)_{(y,f(x);\xi)}\\
    &=\left(\mdiff[f]^{-1}A_{\cotb[Y]}\right)_{(y,f(x);\xi)}
    \tens[]\left(\mpb[f]^{-1}A_{\cotb[X]}
    \right)_{(y,f(x);\xi)}\\
    &=\left(A_{\cotb[Y]}\right)_{(y;f^{\ast}\xi)}
    \tens[]\left(\mpb[f]^{-1}A_{\cotb[X]}
    \right)_{(y;f^{\ast}\xi)}\\
    &=A\\
    &=\left(A_{Y\mathbin{\times_{X}}\cotb[X]}\right)_{(y,f(x);\xi)}
  \end{align*}
  である．
\end{proof}

\(f_1\colon Y\times Y\to X\times Y\)で写像\(f\times \id_Y\)を表し，
\(f_1\colon X\times Y\to X\times X\)で写像\(\id_X\times f\)を表す．
これにより以下の可換図式を得る．
ただし右側の四角はカルテシアンであり\(f_2\)は\(\diag[X]\)に対して横断的である．
\begin{equation}\label{eq:4.4.3}
  \vcenter{\xymatrix@C=36pt@R=36pt{
    Y\times Y
    \ar[r]_-{f_1}
    &
    X\times Y
    \ar@{}[rd]|\square
    %\ar[d]^-{\rotatebox{90}{\(\sim\)}}
    \ar[r]_-{f_2}
    &
    X\times X
    %\ar[d]^-{\rotatebox{90}{\(\sim\)}}
    \\
    \diag[Y]  
    \ar@{^{(}->}[u]^-{}
    \ar[r]^-{{\sim}}
    &
    \diag[]
    \ar@{^{(}->}[u]^-{}
    %\ar[l]^-{\mdiff[f]}
    \ar[r]_-{f}
    &
    \diag[X].
    \ar@{^{(}->}[u]^-{}
  }}
\end{equation}
\begin{proof}
  カルテシアンであることは閉うめこみから従う．

  \(f_2\)が\(\diag[X]\)に対して横断的であることを示す．
  \(\mres[{\mdiff[f_2]}]{(X\times Y)\pb[X\times X]T^{\ast}_{\diag[X]}(X\times X)}\)が単射であればよい．
  %この写像は局所的に\[
  %  (f(y),y;\xi,0)\mapsto (f(y),y;\xi,0)
  %\]とかけるので完了．
  まず，\(\mdiff[f_2]\)は
  \[    
    \vcenter{\xymatrix@C=26pt@R=10pt{
      (X\times Y)\pb[X\times X]T^{\ast}(X\times X)
      %\ar[d]_-{\tilde{f'}}
      \ar[r]^-{\mdiff[f_2]}
      %\ar@{}[dr]|\square
      &
      T^{\ast}(X\times Y)
      %\ar[d]^-{f}
      \\
      \left(\left(f(y),y\right),(f(y),f(y);\xi_1,\xi_2)\right)
      \ar@{}[u]|{\rotatebox{90}{\(\in\)}}
      \ar@{{|}->}[r]_-{}
      %\ar@{}[dr]|\square
      &
      \ar@{}[u]|{\rotatebox{90}{\(\in\)}}
      \left(f(y),y;f_{2}^{\ast}(\xi_1,\xi_2)\right)
    }}
  \]であり，
  \[
    f_2^{\ast}(\xi_1,\xi_2)
    =
    (\id_X\times f)^{\ast}(\xi_1,\xi_2)
    =(\xi_1,f^{\ast}\xi_2)
  \]であるから，\(T^{\ast}_{\diag[X]}(X\times X)\)の点が局所的に
  \((x,x;\xi,-\xi)\)とかけることに注意すると，
  \[
    \mres[{\mdiff[f_2]}]{(X\times Y)\pb[X\times X]T^{\ast}_{\diag[X]}(X\times X)}
    \left(\left(f(y),y\right),(f(y),f(y);\xi,-\xi)\right)
    =(f(y),y;\xi,-f^{\ast}\xi)
  \]
  とかける．\((f(y),y;\xi,-\xi)\)に対し，\((f(y),y;\xi,-f^{\ast}\xi)=(f(y),y;0,0)\)となるとすると，
  \(\xi=0,-f^{\ast}\xi=0\)より，\(\xi=-\xi=0\)となるので単射である．
\end{proof}
次の記号を定める．
\begin{itemize}
  \item \(q_j\)：\(X\times X\)上の第\(j\)射影
  \item \(\tilde{q}_{j}\)：\(X\times Y\)上の第\(j\)射影
  \item \(q'_j\)：\(Y\times Y\)上の第\(j\)射影
\end{itemize}
\begin{DFN}
  \(G\in\Dompb(Y)\), \(F\in\Dompb(X)\)に対し，次のようにおく．
  \begin{enumerate}[(i)]
    \item \(\muhom(G\to F)\coloneqq\mu_{\diag[]}\RHOM_{A_{X\times Y}}(\tilde{q}_2^{-1}G,\tilde{q}_1^{!}F)\)
    \item \(\muhom(F\leftarrow G)\coloneqq\left(\mu_{\diag[]}\RHOM_{A_{X\times Y}}(\tilde{q}_2^{-1}G,\tilde{q}_1^{!}F)\right)^{a}\)
    \item \(Y\to X\)で\(f\)が恒等射のとき，\[
      \muhom(G,F)
      \coloneqq\muhom(G\to F)
      =\mu_{\diag[X]}\RHOM_{A_{X\times X}}({q}_2^{-1}G,{q}_1^{!}F)
    \]
  \end{enumerate}
\end{DFN}
(ii)の右辺の\((\blk)^{a}\)は\(Y\tpb[X]\cotb[X]\)上の
対蹠写像による逆像を表す．

\(\pi\)で\(T^{\ast}_{\diag[]}(X\times Y)\)から\(\diag[]\cong Y\)への射影を表す．

\begin{leftbar}
\begin{PRP}
  次の標準的な同型が存在する．
  \begin{enumerate}[(i)]
    \item 一般に
    %\vspace{-2\baselineskip}
    \begin{align*}
      %\vspace{-2\baselineskip}
      \Rder{\pi_{\ast}}\muhom(G\to F)&\cong\RHOM(G,f^{!}F),\\
      %\vspace{-2\baselineskip}
      \Rder{\pi_{\ast}}\muhom(F\leftarrow G)&\cong\RHOM(f^{-1}F,G).
    \end{align*}
    \(Y=X\)で\(f\)が恒等射の場合はとくに\[
      \Rder{\pi_{\ast}}\muhom(G, F)\cong\RHOM(G,F)
    \]
    となる．
    \item \(G\)がコホモロジー構成可能層と仮定する．このとき次が成り立つ．\[
      \Rder{\pi_{!}}\muhom(G\to F)
      \cong
      \RHOM(G,A_Y)\lltens[]f^{-1}F\tens[]\omega_{Y/X}.
    \]
    \(F\)がコホモロジー構成可能層のときは
    \[
      \Rder{\pi_{!}}\muhom(F\leftarrow G)
      \cong
      f^{-1}\RHOM(f^{-1}F,A_X)\lltens[]G.
    \]が成り立つ．
    \(Y=X\)で\(f\)が恒等射の場合はとくに\[
      \Rder{\pi_{!}}\muhom(G, F)\cong\RHOM(G,A_X)\lltens[]F
    \]
    となる．
  \end{enumerate}
\end{PRP}
\end{leftbar}
\begin{proof}
  (i) 
  定理4.3.2から\begin{align*}
    \Rder{\pi_{\ast}}\mu_{\diag[]}\RHOM(\tilde{q}_2^{-1}G,\tilde{q}_1^{!}F)
    &\underset{(1)}{\cong}
    \Rder{{\tilde{q_2}}_{\ast}}\RG_{\diag[]}\RHOM(\tilde{q}_2^{-1}G,\tilde{q}_1^{!}F)\\
    &\underset{(2)}{\cong}
    \Rder{{\tilde{q_2}}_{\ast}}\RHOM(\tilde{q}_2^{-1}G,\RG_{\diag[]}\tilde{q}_1^{!}F)\\
    &\underset{(3)}{\cong}
    \RHOM(G,\Rder{{\tilde{q_2}}_{\ast}}\RG_{\diag[]}\tilde{q}_1^{!}F)\\
    &\underset{(4)}{\cong}
    \RHOM(G,f^{!}F)
  \end{align*}
  である．ただし，それぞれの同型の根拠は次のとおり．
  \begin{enumerate}[(1)]
    \item 定理4.3.2 (iv) の同型\(
      \mres[\mu_M(F)]{M}\cong i^{!}F
    \)と同一視\(Y\cong\diag[]\)から従う．
    詳しくは以下のようになる．
    \(i\colon \diag[]\hookrightarrow X\times Y\)を閉うめこみ，
    \(\varphi\colon\diag[]\to Y\)を射影，
    \(\psi\colon Y\to \diag[]\)をその逆写像とする．
    このとき，\begin{align*}
      \Rder{\pi_{\ast}}\mu_{\diag[]}
      &\cong
      i^{!}&&\text{（ここに定理4.3.2）}\\
      &\cong
      \Rder{\varphi_{\ast}}i^{!}&&\text{（同一視からくる同型）}\\
      &\cong
      \Rder{\tilde{q_2}_{\ast}}\Rder{i_{\ast}}i^{!}&&\text{（図式の可換性）}\\
      &\cong
      \Rder{\tilde{q_2}_{\ast}}\RG_{\diag[]}&&\text{（\(\RG_{\diag[]}\)の書き換え）}
    \end{align*}
    である．
    \item \(\RG_{\diag[]}\)に関する公式(2.6.9)
    \item \(\tilde{q}_2\)に関する順像逆像随伴
    \item これも同一視からくる．詳しくかくと，
    図式\[\vcenter{\xymatrix{
      &
      X\times Y
      \ar[ld]_-{\tilde{q}_1}
      \ar[rd]^-{\tilde{q}_2}&\\
      X&
      \diag[]
      \ar[l]^-{f}
      \ar[u]_-{i}
      \ar[r]_-{\varphi}
      &Y
    }}\]を用いて，
    \begin{align*}
      f^{!}F
      &\cong i^{!}\tilde{q_1}^{!}F&&\text{（図式の可換性）}\\
      &\cong \Rder{\varphi_{\ast}}i^{!}\tilde{q_1}^{!}F&&\text{（同一視からくる同型）}\\
      &\cong \Rder{\tilde{q_2}_{\ast}}\Rder{i_{\ast}}i^{!}\tilde{q_1}^{!}F&&\text{（図式の可換性）}\\
      &\cong \Rder{\tilde{q_2}_{\ast}}\RG_{\diag[]}\tilde{q_1}^{!}F&&\text{（\(\RG_{\diag[]}\)の書き換え）}
    \end{align*}
    のようになる．
  \end{enumerate}
  2つ目の式は，\(\pi=\pi\circ a\)に注意すれば同様に\begin{align*}
    \Rder{\pi_{\ast}}(\mu_{\diag[]}\RHOM(\tilde{q}_1^{-1}F,\tilde{q}_{2}^{!}G))^{a}
    &\underset{}{\cong}
    \Rder{\pi_{\ast}}\mu_{\diag[]}\RHOM(\tilde{q}_1^{-1}F,\tilde{q}_{2}^{!}G)\\
    &\underset{}{\cong}
    \Rder{{\tilde{q_2}}_{\ast}}\RG_{\diag[]}\RHOM(\tilde{q}_{1}^{-1}F,\tilde{q}_2^{!}G)\\
    &\underset{(2.6.9)}{\cong}
    \Rder{{\tilde{q_2}}_{\ast}}\RHOM((\tilde{q}_1^{-1}F)_{\diag[]},\tilde{q}_2^{!}G)\\
    &\underset{(V.D.)}{\cong}
    \RHOM(\Rder{{\tilde{q_2}}_{!}}(\tilde{q}_1^{-1}F)_{\diag[]},G)\\
    &\underset{}{\cong}
    \RHOM(f^{-1}F,G)
  \end{align*}
  である．

  \(Y=X\)で\(f=\id_X\)のときは，
  \begin{align*}
    \Rder{\pi_{\ast}}\muhom(G,F)
    &=\Rder{\pi_{\ast}}\muhom(G\to F)\\
    &\cong\RHOM(G,\id_{X}^{!}F)\\
    &=\RHOM(G,F)
  \end{align*}である．

  (ii) 
  命題3.4.4.\ref{prp3.4.4}より，
  \begin{align*}
    \Rder{\pi_{!}}\mu_{\diag}\RHOM(\tilde{q}_2^{-1}G,\tilde{q}_1^{!}F)
    &\cong
    \Rder{\pi_{!}}\mu_{\diag}(\rmD{G}\letens F)\\
    &\cong
    i^{-1}(\rmD{G}\letens F)\tens[] \omega_{\diag/{X\times Y}}\\
    &\cong
    i^{-1}\tilde{q}_2^{-1}\rmD{G}\tens i^{-1}\tilde{q}_1^{-1}F\tens[] \omega_{\diag/{X\times Y}}\\
    &\cong
    \varphi^{-1}\rmD{G}\tens f^{-1}F\tens[] \omega_{\diag/{X\times Y}}\\
    &\cong
    \rmD{G}\tens f^{-1}F\tens[] \omega_{\diag/{X\times Y}}
  \end{align*}
\end{proof}

関手\(\muhom\)は佐藤超局所化の関手を一般化したものである．

\begin{leftbar}
\begin{PRP}
  \(Y\)を\(X\)の閉部分多様体とし，
  \(j\)をうめこみ\(\conm[Y]{X}\to \cotb[X]\)とする．
  \(F\in\Dompb(X)\)とすると，次の同型が存在する．
  \[
    \muhom(A_Y,F)\cong j_{\ast}\mu_{Y}F.
  \]
\end{PRP}
\end{leftbar}
\begin{proof}
  \(f\)をうめこみ\(Y\hookrightarrow X\)とする．このとき図式\[
  \vcenter{\xymatrix{
      X\times Y
      %\ar[ld]_-{\tilde{q}_1}
      \ar@{^{(}->}[r]^-{f_2}
      \ar[d]^-{\tilde{q}_1}
      &
      X\times X
      \ar[ld]^-{q_1}
      \\
      X
      %\ar[l]^-{f}
      %\ar[u]_-{i}
      %&Y
  }}\]を考えて，
  \begin{align*}
    \mu_{\diag[X]}\RHOM(q_2^{-1}A_{Y},q_1^!F)
    &\cong
    \mu_{\diag[X]}\RG_{X\times Y}(q_1^!F)\\
    &\cong
    \mu_{\diag[X]}(\Rder{{f_2}_{\ast}}{f_2}^{!}q_1^!F)\\
    &\cong
    \mu_{\diag[X]}(\Rder{{f_2}_{\ast}}\tilde{q}_{1}^{!}F)
  \end{align*}
  である．
  \(\conm[{\diag[]}]{(X\times Y)}\)を
  \(
    Y\tpb[X]\conm[{\diag[X]}]{(X\times X)}
  \)と同一視する\footnote{
    同型\(\cotb[X]\cong\conm[{\diag[X]}]{(X\times X)}\)に引き戻し\(Y\tpb[X](\blk)\)を適用すると
    \(
      Y\tpb[X]\cotb[X]
      \cong 
      Y\tpb[X]\conm[{\diag[X]}]{(X\times X)}
    \)となる．これと\(
      Y\tpb[X]\cotb[X]
      \cong 
      \conm[{\diag[]}]{(X\times Y)}
    \)から同型を得る．
  }．座標でかくと\[    
    \vcenter{\xymatrix@C=26pt@R=10pt{
      \conm[{\diag[]}]{(X\times Y)}
      %\ar[d]_-{\tilde{f'}}
      \ar[r]^-{\sim}
      %\ar@{}[dr]|\square
      &
      Y\dpb[X]\conm[{\diag[X]}]{(X\times X)}
      %\ar[d]^-{f}
      \\
      \left(f(y),y;\xi,-f^{\ast}\xi\right)
      \ar@{}[u]|{\rotatebox{90}{\(\in\)}}
      \ar@{{|}->}[r]_-{}
      %\ar@{}[dr]|\square
      &
      \ar@{}[u]|{\rotatebox{90}{\(\in\)}}
      \left(y,f(y),f(y);\xi,-\xi\right)
    }}
  \]である．
  \begin{CMT}
    \(\mres[f_2]{\diag}\colon\diag[]\to\diag[X]\)の基本図式
    \[
      \vcenter{\xymatrix@C=36pt@R=36pt{
      T^{\ast}_{\diag[]}(X\times Y)
      \ar[dr]^-{}
      &
      \diag[]\dpb[{\diag[X]}]T^{\ast}_{\diag[X]}(X\times X)
      \ar[d]^-{}
      \ar[l]_-{\mdiff[{f_2}_{\diag[]}]}
      \ar[r]^-{\mpb[{f_2}_{\diag[]}]}
      &
      T^{\ast}_{\diag[X]}(X\times X)
      \ar[d]^-{}
      \\
      %\cotb[Y]
      &
      \diag[]
      %\ar[l]^-{}
      \ar[r]_-{\mres[{f_2}]{\diag[]}}
      &
      \diag[X]
    }}\]
    を考える．
    これは同一視\(\varphi\colon \diag[]\to Y\), 
    \(\psi\colon \diag[X]\to X\), 
    \(\conm[{\diag[X]}]{(X\times X)}\simar \cotb[X]\)のもとで
    \[
      \vcenter{\xymatrix@C=36pt@R=36pt{
      T^{\ast}_{\diag[]}(X\times Y)
      \ar[dr]^-{}
      &
      Y\dpb[X]T^{\ast}_{\diag[X]}(X\times X)
      \ar[d]^-{}
      \ar[l]_-{\mdiff[f]}
      \ar[r]^-{\mpb[f]}
      &
      \cotb[X]
      \ar[d]^-{}
      \\
      %\cotb[Y]
      &
      Y
      %\ar[l]^-{}
      \ar[r]_-{f}
      &
      X
    }}\]
    となる．
    したがって，上の同一視は\(
      \mdiff[f]=\mdiff[{f_2}_{\diag[]}]
    \)によるものである．
  \end{CMT}

  \(f_2\)は\(\diag[X]\)に対して横断的なので特に鮮明である．
  また\(f_2\)は閉うめこみなので固有的である．
  \(f_2^{-1}(\diag[X])=\diag[]\)も明らか．
  したがって命題\ref{prp4.3.4}から次が成り立つ．
  \begin{align*}
    \mu_{\diag[X]}(\Rder{{f_2}_{\ast}}\tilde{q}_{1}^{!}F)
    &\cong
    \Rder{{\mpb[{f_2}_{\diag[]}]}_{\ast}}
    {\mdiff[{f_2}_{\diag[]}]^{-1}}
    \mu_{\diag[]}(\tilde{q}_{1}^{!}F)\\
    &\cong
    \Rder{{\mpb[{f_2}_{\diag[]}]}_{\ast}}
    \mu_{\diag[]}(\tilde{q}_{1}^{!}F).
    %\quad
    &&\text{（同型\(\mdiff[{{f_2}_{\diag[]}}]\)による同一視）}
  \end{align*}
  さらに，命題\ref{prp4.3.5}
  より，
  \begin{align*}
    \Rder{{\mpb[{f_2}_{\diag[]}]}_{\ast}}\mu_{\diag[]}(\tilde{q}_{1}^{!}F)  
    &\cong
    \Rder{{\mpb[{f_2}_{\diag[]}]}_{\ast}}
    \Rder{{\mdiff[({{\widetilde{q_{1}}}_{\diag[]}})]}_{\ast}}
    \mpb[({\widetilde{q_{1}}_{\diag[]}})]^{!}
    \mu_{Y}(F)\\
    &\cong
    \Rder{{\mpb[{f_2}_{\diag[]}]}_{\ast}}
    \Rder{{\mdiff[({{\widetilde{q_{1}}}_{\diag[]}})]}_{\ast}}
    \mu_{Y}(F)
    %\quad
    &&\text{（同型\(\mpb[({\widetilde{q_{1}}_{\diag[]}})]\)による同一視）}\\
    &\cong
    j_{\ast}\mu_{Y}(F)        
    %\quad
    &&\text{（図式の可換性）}
  \end{align*}
  となる．
\end{proof}









% ======================================================
% appendix
% ======================================================
\appendix
\section{ノート中で用いる主張}
\begin{leftbar}
\begin{PRP}[{\cite[Prop.3.4.4.]{KS90}}]\label{prp3.4.4}
  \(X\)と\(Y\)を有限c柔軟次元局所コンパクト（ハウスドルフ）空間とし，
  \(q_1\)と\(q_2\)をそれぞれ\(X\times Y\)から\(X\)と\(Y\)への射影とする．
  \(F\in\Dompb(A_X)\), \(G\in\Dompb(A_Y)\)とし，\(F\)はコホモロジー構成可能層とする．
  このとき，
  \[
    \vD{F}\letens G\to \RHOM(q_1^{-1}F,q_2^{!}G)
  \]は同型である．
  \(X\)が\(C^0\)多様体ならば，\[
    \vD'{F}\letens G\to \RHOM(q_1^{-1}F,q_2^{-1}G)
  \]も同型である．
\end{PRP}
\end{leftbar}

\begin{leftbar}
\begin{PRP}[{\cite[Prop.4.3.4]{KS90}}]\label{prp4.3.4}
  \(G\in\Dompb(Y)\)とする．
  このとき，標準的な射による可換図式    
  \[    
      \vcenter{\xymatrix@C=36pt@R=36pt{
      \Rder{{\mpb[f_{N}]}_{!}}{\mdiff[f_{N}]^{-1}}\mu_{N}(G)
      \ar[d]
      \ar[r]
      %\ar@{}[dr]|\square
      &
      \mu_{M}(\Rder{f_{!}}G)
      %\ar[r]
      \ar[d]
      \\
      \Rder{{\mpb[f_{N}]}_{\ast}}({\mdiff[f_{N}]^{!}}\mu_{N}(G)\tens[]\omega_{Y/X}\tens[]\omega_{N/M}^{\otimes{-1}})
      %\ar@{}[dr]|\square
      &
      \mu_{M}(\Rder{f_{\ast}}G)
      \ar[l]
      }}
  \]が存在する．
  さらに，\(\supp(G)\to X\)と\(
      C_N(\supp(G))\to T_{M}X
  \)が固有かつ\(\supp(G)\cap f^{-1}(M)\subset N\)ならば，
  以上の射はすべて同型である．
  
  とくに\(f^{-1}(M)=N\)かつ，
  \(f\)が\(M\)に関して鮮明\footnotemark[2] (clean) かつ\(\supp(G)\)上固有のとき，
  以上の射はすべて同型である．
\end{PRP}
\end{leftbar}
\footnotetext[2]{試験的な訳語}
\begin{leftbar}
\begin{PRP}[{\cite[Prop.4.3.5]{KS90}}]\label{prp4.3.5}
  \(F\in\Dompb(X)\)とする．
  \begin{enumerate}[(i)]
      \item 標準的な射で次の図式を可換にするものが存在する．
      \[\vcenter{\xymatrix@C=56pt@R=46pt{
          \Rder{{\mdiff[f_N]}_{!}}(\omega_{N/M}\tens[]{\mpb[f_N]^{-1}}\mu_{M}(F))
          \ar[d]^{}
          \ar[r]_{}
          %\ar@{}[dr]|\square
          &
          \mu_{N}(\omega_{Y/X}\tens[]{f^{-1}}F)
          %\ar[r]
          \ar[d]
          \\
          \Rder{{\mdiff[f_N]}_{\ast}}{\mpb[f_N]^{!}}\mu_{M}(F)
          %\ar@{}[dr]|\square
          &
          \mu_{N}({f^{!}}F)
          \ar[l]_{\beta}
      }}\]
      ここで縦向きの矢印は(3.1.6)から導かれるものである．
      \item \(f\colon Y\to X\)と\(\mres[f]{N}\colon N\to M\)が平滑ならばこれらの射はすべて同型となる．    
      \item \(f\)が\(M\)に横断的かつ\(f^{-1}(M)=N\)ならば，
      自然な射\[
          \Rder{{\mdiff[f_N]}_{\ast}}\mpb[f_N]^{-1}\mu_{M}(F)
          \to \mu_{N}(f^{-1}F)
      \]が存在する．
  \end{enumerate}
\end{PRP}
\end{leftbar}















%===============================================
% 参考文献スペース
%===============================================
\begin{thebibliography}{20} 
  \bibitem[dR]{dR} 
  de Rham, \textit{Vari\'et\'es diff\'erentiables}, Hermann, 1953.
  和訳：ド・ラーム, 微分多様体, 東京図書, 1974.
  \bibitem[GKS12]{GKS12} St\'ephane Guillermou, 
  Masaki Kashiwara, Pierre Schapira, 
  \textit{Sheaf Quantization of Hamiltonian Isotopies and Applications to Nondisplaceability Problems}, 
  Duke. Math. J. 161, No. 2, pp.201--245, 2012.  
  \bibitem[GP74]{GP74} Victor Guillemin, Alan Pollack, 
    \textit{Differential Topology}, 
    Prentice-Hall, 1974.
  \bibitem[KS90]{KS90} Masaki Kashiwara, Pierre Schapira, 
    \textit{Sheaves on Manifolds}, 
    Grundlehren der Mathematischen Wissenschaften, 292, Springer, 1990.
  \bibitem[Hat89]{Hat89} 服部晶夫, 多様体 増補版, 岩波全書 288, 岩波書店, 1989.
  %\bibitem[Og02]{Og02} 小木曽啓示, 代数曲線論, 朝倉書店, 2022.
  \bibitem[Hor63]{Hor63} Lars H\"ormander, 
  \textit{Linear Pertial Differential Operators}, Fourth Printing, 
  Grundlehren der Mathematischen Wissenschaften, 116, Springer, 1963.
  \bibitem[Hor71]{Hor71} Lars H\"ormander, 
  \textit{Fourier Integral Operators I}, 
  Acta Math. 127, 79--183, (1971).
  \bibitem[Hor89]{Hor89} Lars H\"ormander, 
  \textit{The Analysis of Linear Pertial Differential Operators I}, Second Edition,
  Grundlehren der Mathematischen Wissenschaften, 256, Springer, 1989.
  \bibitem[Hor85]{Hor85} Lars H\"ormander, 
  \textit{The Analysis of Linear Pertial Differential Operators III}, 
  Grundlehren der Mathematischen Wissenschaften, 274, Springer, 1985.
  \bibitem[Sch85]{Sch85} Pierre Schapira, 
    \textit{Microdifferential Systems in the Complex Domain}, 
    Grundlehren der Mathematischen Wissenschaften, 269, 
    Springer, 1985.
    \bibitem[ST92]{ST92} Pierre Schapira, Nobuyuki Tose, 
    \textit{Morse Inequalities for \(\rr\)-constructible Sheaves}, 
    Adv. Math. 93, pp.1--8, 1992. 
    \bibitem[Zwo12]{Zwo12} Maciej Zworski,
    \textit{Semiclsaaical Analysis}, 
    Graduate Studies in Mathematics, 138, 
    American Mathematical Society, 
    2012. 
\end{thebibliography}

%===============================================




\end{document}
